\documentclass[UTF8,a4paper,12pt]{ctexart}
\usepackage[margin=2.5cm]{geometry}
\usepackage{amsmath,booktabs,hyperref,xcolor}
\hypersetup{colorlinks=true,linkcolor=blue,urlcolor=blue}
\title{课程项目报告\\\large Coursework Report Starter}
\author{姓名 Your Name\\学号 Student ID: 20260001}
\date{\today}
\begin{document}
\maketitle
\begin{center}
课程名称 Course: Data Analysis \qquad 指导教师 Advisor: Name
\end{center}

\begin{abstract}
简要说明问题、方法和主要结论。 This abstract should state the question,
method, evidence, and conclusion in a form that stands alone.
\end{abstract}

\section{问题提出 Problem statement}
描述背景、研究问题和成功标准。 Explain the context, research question, and
what would count as a useful answer.

\section{方法与实现 Method and implementation}
记录数据来源、预处理步骤和实现选择。 For a reproducible project, record
versions, parameters, and assumptions alongside the narrative.
\begin{equation}
  \bar{x}=\frac1n\sum_{i=1}^{n}x_i,\qquad
  s^2=\frac1{n-1}\sum_{i=1}^{n}(x_i-\bar{x})^2.
\end{equation}

\clearpage
\section{实验结果 Results}
\begin{table}[ht]
\centering
\caption{示例实验结果 / Example evaluation}
\begin{tabular}{@{}lrrr@{}}
\toprule
方案 Configuration & 指标 Metric & 误差 Error & 时间 Time (s)\\\midrule
基线 Baseline & 0.742 & 0.031 & 8.4\\
改进 Improved & 0.816 & 0.024 & 9.7\\
\bottomrule
\end{tabular}
\end{table}

\section{分析与讨论 Analysis and discussion}
解释结果与假设的关系，比较替代方案，并讨论不确定性。 Separate observed
facts from interpretation and record limitations that could change the result.

\section{总结 Conclusion}
用简洁的条目回答研究问题，并列出下一步工作。 Summarise the evidence in
two or three concrete points and identify the most valuable follow-up.

\begin{thebibliography}{9}
\bibitem{montgomery} D. C. Montgomery, \emph{Design and Analysis of
  Experiments}, Wiley, 2017.
\end{thebibliography}
\end{document}
